% ============================================================================
% Fichier  : partie6/P06_C32_cas_pratique_zknr.tex
% Titre    : Cas Pratique -- Audit Formel d'un Protocole Multi-Couches
% Auteur   : MINKA MI NGUIDJOÏ Thierry Emmanuel
% Version  : 2.0 -- Juin 2026
% ============================================================================

\chapter{Cas Pratique. Audit Formel d'un Protocole Multi-Couches}
\label{chap:cas-zknr}

\epigraph{%
    « In theory, theory and practice are the same. In practice, they are
    not. The true test of any protocol is in its implementation. »%
}{attribué à Albert Einstein, formulation populaire}

\niveauChapitre{Recherche}{%
    Maîtrise complète de la Partie~VI requise (Ch.~\ref{chap:CRO}
    à~\ref{chap:form-meth}). Ce chapitre applique l'ensemble des
    outils théoriques accumulés à un cas concret.%
}

\begin{boxobjectifs}
\begin{itemize}
    \item Appliquer intégralement la méthodologie d'audit en cinq
          étapes (Ch.~\ref{chap:form-meth}) à un cas concret~:
          l'analyse cryptanalytique de la signature BLS.
    \item Comprendre la vulnérabilité spécifique de BLS face à
          l'algorithme de Shor, et pourquoi cette vulnérabilité
          motive son couplage architectural avec une primitive
          post-quantique.
    \item Maîtriser une checklist opérationnelle d'analyse face à
          une preuve cryptographique multi-couches rencontrée en
          investigation.
    \item Conduire vous-même, à titre d'exercice guidé, les quatre
          premières étapes de l'audit formel d'une architecture
          multi-couches décrite publiquement, en respectant
          scrupuleusement la distinction entre ce qui peut être
          établi à partir d'une description architecturale et ce
          qui nécessiterait un accès aux spécifications complètes.
    \item Synthétiser, sur cet exemple, l'ensemble des outils de
          la Partie~VI~: Trilemme CRO formalisé, méthodologie de
          calcul des primitives, taxonomie cryptanalytique,
          vérification formelle.
\end{itemize}
\end{boxobjectifs}

\bigskip

% ============================================================================
\section*{De la théorie à la pratique~: ce que ce chapitre démontre}
% ============================================================================

Les Chapitres~\ref{chap:CRO} à~\ref{chap:form-meth} ont construit,
brique par brique, l'ensemble des outils nécessaires à l'analyse
rigoureuse d'un système cryptographique~: théorie de l'information,
définitions formelles CLO, taxonomie cryptanalytique, méthodologie
d'audit. Ce chapitre les met en \oe uvre sur un cas réel et entièrement
vérifiable, la signature \textbf{BLS} (Boneh-Lynn-Shacham), une
primitive cryptographique authentique, publiée, abondamment
documentée dans la littérature scientifique, avant de vous guider
dans l'application de cette même méthodologie à une architecture
multi-couches plus ambitieuse, à titre d'exercice de synthèse.

\separateur

% ============================================================================
\section{Étude de Cas Authentique~: la Signature BLS}
\label{sec:etude-bls}
% ============================================================================

\subsection{Fonctionnement et propriétés}

\begin{boxdef}
\textbf{Signature BLS}\index{BLS, signature}~(Boneh, Lynn, Shacham,
2001)~: schéma de signature numérique fondé sur les couplages
bilinéaires (\textit{pairings}) sur courbes elliptiques, typiquement
instancié sur la courbe \texttt{BLS12-381}. Sa propriété distinctive~:
les signatures sont \textbf{courtes} et
\textbf{agrégeables}, plusieurs signatures de plusieurs signataires
sur plusieurs messages peuvent être combinées en une seule signature
compacte, vérifiable efficacement.
\end{boxdef}

Cette propriété d'agrégation rend BLS particulièrement adapté aux
\textbf{systèmes à seuil} (\textit{threshold signatures})~: un
ensemble de signataires peut produire collectivement une signature
unique, sans qu'aucun signataire individuel ne puisse signer seul, pertinent pour des contextes institutionnels où une décision
doit recueillir l'accord de plusieurs parties avant de produire
une attestation juridiquement valable (par exemple, une décision
collégiale en investigation, Ch.~6~(UE~SEC4031), Quality Gates).

\subsection{Cryptanalyse classique et quantique}

\begin{tableaucours}{p{2.5cm}p{4.0cm}p{4.5cm}}{%
    \tableauHeader{Horizon} &
    \tableauHeader{Fondement de sécurité} &
    \tableauHeader{État de la cryptanalyse}
}
\textbf{Classique} & Difficulté du problème calculatoire de
    Diffie-Hellman (CDH) sur courbes appariées &
    Aucune attaque efficace connue sur BLS12-381 à la date de
    rédaction de ce cours, sécurité considérée comme solide
    face aux capacités de calcul classiques actuelles \\
\textbf{Quantique} & \textit{Le même problème CDH, mais résoluble
    par Shor} & \textbf{Extrêmement vulnérable}~: l'algorithme de
    Shor (Ch.~\ref{chap:quantique}) résout le problème CDH en
    temps polynomial sur un ordinateur quantique suffisamment
    puissant, réduisant la sécurité de la signature
    \textbf{à néant}, pas une dégradation progressive, une
    rupture complète \\
\end{tableaucours}
\vspace{-0.8em}
\begin{center}{\small\itshape Tableau~32.1. Cryptanalyse de BLS,
horizons classique et quantique}\end{center}

\begin{boxalert}
\textbf{BLS et le Théorème 1 du Chapitre~\ref{chap:quantique}~:
une application directe}

Cette vulnérabilité de BLS face à Shor n'est pas une particularité
isolée~: elle illustre exactement le Tableau~26.1 du
Ch.~\ref{chap:quantique}, qui classe \textbf{toute} signature fondée
sur le logarithme discret (dont BLS, comme ECDSA) parmi les
primitives vulnérables à Shor. La distinction
hash~\textit{vs.}~signature, déjà établie au Ch.~\ref{chap:quantique},
s'applique ici sans modification~: BLS n'a aucune protection
structurelle particulière contre cette menace, c'est une
signature asymétrique classique, avec les mêmes faiblesses que
toute primitive de sa catégorie.
\end{boxalert}

\subsection{Positionnement qualitatif dans le cadre CRO}

\begin{boiteRemarque}{Rigueur méthodologique~: positionnement
qualitatif, pas un calcul publié}
Conformément à la discipline déjà établie au
Ch.~\ref{chap:prim-crypt} (Section~6, Tableau~28.6), BLS \textbf{ne
dispose pas d'un calcul CLO publié} avec décomposition bit par bit
vérifiable. Le positionnement qui suit est donc qualitatif,
raisonné par analogie avec ECDSA (calcul publié,
Ch.~\ref{chap:prim-crypt}, Section~4), une primitive de
structure comparable (signature asymétrique, transparente par
construction).
\end{boiteRemarque}

\begin{tableaucours}{p{2.5cm}p{3.0cm}p{5.0cm}}{%
    \tableauHeader{Dimension} &
    \tableauHeader{Positionnement attendu} &
    \tableauHeader{Raisonnement par analogie}
}
$\mathsf{Conf}$ & Faible &
    Comme ECDSA (Ch.~\ref{chap:prim-crypt}, $\mathsf{Conf}\approx 0{,}31$)~:
    une signature, par construction, ne masque pas le message
    signé, seule la clé privée reste secrète \\
$\mathsf{Rel}$ & Élevé, \textbf{aujourd'hui} &
    Vérification déterministe~; la falsification exige de
    résoudre le problème CDH, considéré difficile classiquement
    (analogue à la difficulté du logarithme discret elliptique
    pour ECDSA, $\mathsf{Rel}\approx 0{,}98$) \\
$\mathsf{Opp}$ & Probablement élevé, \textbf{aujourd'hui} &
    Jurisprudence émergente autour des signatures à courbes
    elliptiques en général~; cependant moins établie que celle
    d'ECDSA, plus ancien et plus largement déployé dans les
    infrastructures PKI (Ch.~11~(UE~SEC4031)) \\
\textbf{Horizon post-quantique} & \textbf{Effondrement de
    $\mathsf{Rel}$} & Une fois Shor opérationnel à l'échelle requise,
    $\mathsf{Rel}(\text{BLS})$ s'effondre, une signature falsifiable
    n'offre plus aucune fiabilité, indépendamment de son
    opposabilité historique \\
\end{tableaucours}
\vspace{-0.8em}
\begin{center}{\small\itshape Tableau~32.2. Positionnement CRO
qualitatif de BLS (raisonnement par analogie, non calculé)}\end{center}

Ce déséquilibre temporel, solide aujourd'hui, structurellement
condamné demain, justifie, sur le plan architectural, le principe
de \textbf{défense en profondeur} déjà établi au
Ch.~\ref{chap:cryptanal}~: coupler une primitive performante mais
vulnérable au quantique (BLS) avec une primitive post-quantique
plus lourde mais résistante (par exemple, Dilithium/ML-DSA, FIPS~204,
Ch.~\ref{chap:quantique}), de sorte que la sécurité globale du
système hérite de la résistance du \textbf{plus robuste} des deux
mécanismes face à chaque horizon temporel.

% ============================================================================
\section{Recommandations Opérationnelles pour l'Investigateur}
\label{sec:recommandations-investigateur}
% ============================================================================

\subsection{Face à une signature BLS isolée}

\begin{boxPNC}
\textbf{Protocole d'évaluation d'une signature BLS rencontrée en
investigation}

\textbf{Aujourd'hui}~: une signature BLS valide constitue une preuve
forte d'authenticité et d'intégrité, vérifiable selon les méthodes
standard (Ch.~\ref{chap:rapport}, distinction constatation~:
``la signature BLS vérifie correctement selon l'équation de
vérification publiée'', un fait, indépendant de toute
interprétation).

\textbf{Pour des preuves destinées à un archivage de longue durée}~:
dans une logique \textit{Harvest Now, Decrypt Later}
(Ch.~\ref{chap:quantique}), considérer explicitement qu'une
signature BLS isolée pourrait devenir forgeable dans le futur.
\textbf{Recommandation~:} rechercher systématiquement si une
signature post-quantique complémentaire (Dilithium ou équivalent
FIPS~204) a été archivée en parallèle pour la même pièce, et
documenter son absence le cas échéant comme une limitation du
dossier (sur le modèle des réserves explicites du
Ch.~\ref{chap:rapport}).
\end{boxPNC}

\subsection{Checklist générale d'analyse d'un protocole}

Cette checklist, synthèse opérationnelle de l'ensemble de la
Partie~VI, s'applique à \textbf{tout} protocole cryptographique
rencontré en investigation, qu'il soit simple (une seule primitive)
ou complexe (architecture multi-couches)~:

\begin{enumerate}
    \item[$\square$] Identifier \textbf{toutes} les primitives
          cryptographiques employées (Ch.~\ref{chap:prim-crypt}).
    \item[$\square$] Évaluer leur résistance classique \textbf{et}
          post-quantique pour chacune (Ch.~\ref{chap:quantique},
          Tableau~26.1).
    \item[$\square$] Cartographier les flux de messages et les
          états persistants (Ch.~\ref{chap:form-meth},
          Étape~1~: Compréhension).
    \item[$\square$] Rechercher les vulnérabilités de composition
          et de rejeu entre primitives (Ch.~\ref{chap:cryptanal},
          taxonomie ``contre le protocole''~; Ch.~\ref{chap:form-meth},
          Étape~3).
    \item[$\square$] Vérifier la présence de mécanismes explicites
          de gestion de clés (rotation, révocation), un défaut
          fréquent et souvent sous-estimé dans les architectures
          multi-couches (cf.~Section~\ref{sec:exercice-guide}).
    \item[$\square$] Consulter les preuves formelles disponibles,
          ou, à défaut, modéliser le protocole dans Tamarin ou
          ProVerif (Ch.~\ref{chap:form-meth}).
    \item[$\square$] Distinguer systématiquement, dans toute
          conclusion, ce qui est \textbf{établi et vérifiable}
          de ce qui relève d'une \textbf{hypothèse} ou d'une
          \textbf{question ouverte} (principe transversal de
          la Partie~VI, rappelé à chaque chapitre).
\end{enumerate}

% ============================================================================
\section{Exercice Guidé~: Audit d'une Architecture Multi-Couches}
\label{sec:exercice-guide}
% ============================================================================

\subsection{Contexte de l'exercice}

Le Ch.~\ref{chap:cryptanal} a mentionné, avec la prudence
méthodologique requise, le projet de recherche \textbf{ZK-NR}
(\textit{Zero-Knowledge Non-Repudiation}) de l'auteur de ce cours~:
une architecture combinant quatre couches, engagements de
Merkle, preuves STARK, signatures à seuil BLS, signatures
post-quantiques Dilithium, pour des scénarios de non-répudiation
à forte criticité.

\begin{boxalert}
\textbf{Ce que cet exercice est, et ce qu'il n'est pas}

Cet exercice \textbf{n'est pas} la présentation de résultats d'audit
établis~: ZK-NR n'existe qu'en prépublication~\cite{minka2025zknr},
sans évaluation par les pairs ni audit indépendant, et ce serait contraire à la rigueur
méthodologique pratiquée depuis le Ch.~\ref{chap:rapport}. Cet
exercice \textbf{est} une mise en situation~: à partir de la seule
description architecturale publique (les quatre couches mentionnées
ci-dessus), vous allez appliquer vous-même les quatre premières
étapes de la méthodologie d'audit (Ch.~\ref{chap:form-meth},
Tableau~31.3), en vous limitant strictement à ce qui peut être
raisonné à partir des propriétés \textbf{publiquement connues} de
chaque primitive composante (Merkle~: Ch.~\ref{chap:for-sys},
intégrité de structures de données~; STARK~: transparent, pas de
\textit{trusted setup}, mentionné au Ch.~\ref{chap:quantique}~;
BLS~: ce chapitre, Section~\ref{sec:etude-bls}~; Dilithium~:
FIPS~204, Ch.~\ref{chap:quantique}).
\end{boxalert}

\subsection{Étape~1. Compréhension (à compléter par vous)}

À partir des quatre couches annoncées, formulez~:
\begin{itemize}
    \item l'objectif de sécurité visé par la combinaison de ces
          quatre primitives (indice~: que permet une preuve à
          divulgation nulle que ne permet pas une simple signature~? Que
          permet une primitive post-quantique que ne permet pas une
          primitive classique~?)~;
    \item le rôle probable de chaque couche selon ses propriétés
          connues (Merkle~: structure d'engagement immuable~;
          STARK~: preuve sans divulgation~; BLS~: signature
          agrégeable à seuil~; Dilithium~: authentification résistante
          au quantique).
\end{itemize}

\subsection{Étape~2. Modélisation (à compléter par vous)}

En vous appuyant sur le Tableau~31.1 (Ch.~\ref{chap:form-meth}),
formulez, en langage naturel d'abord, les lemmes de sécurité que
vous attendriez d'une telle architecture, par exemple~:
``une attestation produite par ce protocole ne peut être reniée
par son émetteur'' (non-répudiation), ``le contenu sous-jacent à
une attestation reste confidentiel sauf divulgation volontaire''
(confidentialité).

\subsection{Étape~3. Analyse manuelle guidée}

Cette étape peut être conduite, en toute rigueur, à partir des
\textbf{seules propriétés publiques} des quatre primitives. Voici
une amorce structurée, à vous de la compléter et de
l'approfondir~:

\begin{tableaucours}{p{3.0cm}p{3.5cm}p{4.0cm}}{%
    \tableauHeader{Aspect} &
    \tableauHeader{Ce que les propriétés publiques
    permettent d'affirmer} &
    \tableauHeader{Ce qui resterait à vérifier}
}
Modularité & Une architecture en couches indépendantes limite,
    en principe, la propagation d'une défaillance d'une couche
    aux autres (principe général de défense en profondeur,
    Ch.~\ref{chap:cryptanal}) & Le degré d'indépendance \textbf{réel}
    entre les couches dans l'implémentation spécifique~;
    une dépendance cachée romprait cette garantie \\
Vulnérabilité BLS & Section~\ref{sec:etude-bls} de ce chapitre
    établit, sur la base de la littérature publique, la
    vulnérabilité de BLS à Shor & Si et comment la couche
    Dilithium compense spécifiquement et complètement cette
    vulnérabilité dans l'architecture précise de ce protocole \\
Transparence STARK & Les STARKs, par construction publique
    documentée, ne requièrent pas de \textit{trusted setup}
    (contrairement aux zk-SNARKs, Ch.~\ref{chap:prim-crypt})
    & La taille effective des preuves générées et leur impact
    sur la performance opérationnelle du protocole \\
Gestion des clés & La checklist générale
    (Section~\ref{sec:recommandations-investigateur}) identifie
    la gestion des clés comme un point de vigilance \textbf{générique}
    pour toute architecture multi-couches & Si des règles formelles
    de rotation et de révocation existent spécifiquement pour ce
    protocole, information non vérifiable sans accès aux
    spécifications complètes \\
\end{tableaucours}
\vspace{-0.8em}
\begin{center}{\small\itshape Tableau~32.3. Amorce d'analyse manuelle,
fondée uniquement sur les propriétés publiques des primitives composantes}\end{center}

\subsection{Étape~4. Vers l'analyse automatisée}

Sur la base de ce que vous avez établi aux étapes précédentes,
esquissez (sans nécessairement l'exécuter) la structure d'un modèle
Tamarin pour une telle architecture~: combien de règles de
réécriture anticipez-vous au minimum (une par couche~? davantage
pour modéliser leurs interactions~?), et quels lemmes prioriseriez-vous
en premier compte tenu des points de vigilance identifiés à
l'Étape~3~?

\begin{boiteRemarque}{La valeur pédagogique de cet exercice}
Que vous ayez accès, un jour, aux spécifications complètes d'un
protocole tel que ZK-NR ou à tout autre protocole multi-couches
rencontré dans votre pratique professionnelle, la démarche que
vous venez de mettre en \oe uvre dans cette section
\textbf{constitue exactement la méthode à appliquer}~:
raisonner à partir de ce qui est publiquement établi, identifier
explicitement les questions qui resteraient ouvertes sans
information supplémentaire, et ne jamais combler ces lacunes par
des affirmations non vérifiables, la compétence centrale de
toute la Partie~VI de ce cours.
\end{boiteRemarque}

\separateur

\begin{boxresume}
\textbf{Ce qu'il faut retenir de ce chapitre~:}
\begin{itemize}
    \item \textbf{BLS}~: signatures courtes et agrégeables, fondées
          sur la difficulté du problème CDH sur courbes appariées~;
          \textbf{solides classiquement, mais structurellement
          condamnées par l'algorithme de Shor}, aucune
          dégradation progressive, une rupture complète.

    \item Le déséquilibre temporel de BLS (robuste aujourd'hui,
          vulnérable demain) \textbf{justifie architecturalement}
          son couplage avec une primitive post-quantique
          (Dilithium/FIPS~204), application directe du principe
          de défense en profondeur (Ch.~\ref{chap:cryptanal}).

    \item \textbf{Recommandation opérationnelle clé}~: pour toute
          signature BLS isolée destinée à un archivage de longue
          durée, vérifier systématiquement l'existence d'une
          signature post-quantique complémentaire archivée en
          parallèle, sinon, documenter cette absence comme
          limitation du dossier.

    \item La \textbf{checklist générale en sept points}
          (Section~\ref{sec:recommandations-investigateur}) est
          directement applicable à tout protocole cryptographique
          rencontré en investigation, simple ou complexe.

    \item Face à une architecture multi-couches \textbf{non
          entièrement documentée} (comme un protocole de recherche
          en cours), la méthode rigoureuse consiste à
          \textbf{raisonner à partir des seules propriétés
          publiques} des primitives composantes, et à
          \textbf{identifier explicitement} ce qui resterait à
          vérifier, jamais à combler les lacunes par des
          affirmations non sourcées.

    \item Cette discipline, distinguer constamment l'établi de
          l'hypothétique, est la compétence unificatrice de
          toute la Partie~VI de ce cours, qu'il s'agisse d'un
          théorème mathématique prouvé (Ch.~\ref{chap:CRO}), d'un
          calcul publié (Ch.~\ref{chap:prim-crypt}), d'un résultat
          de vérification formelle (Ch.~\ref{chap:form-meth}), ou
          d'un protocole en cours de développement (ce chapitre).
\end{itemize}
\end{boxresume}

\bigskip

\begin{boxexo}[title={\ding{45}\quad Exercice~32.1. Cryptanalyse appliquée \niveaubase}]
\begin{enumerate}[(a)]
    \item Expliquez, en vous appuyant sur le Tableau~32.1, pourquoi
          la vulnérabilité de BLS à Shor est qualifiée
          d'``effondrement'' plutôt que de simple
          ``dégradation'', en quoi cela diffère-t-il de l'effet
          de Grover sur AES-256 (Ch.~\ref{chap:quantique})~?
    \item Une institution utilise exclusivement BLS pour signer ses
          archives judiciaires destinées à une conservation de
          50~ans. Identifiez le risque principal et la recommandation
          de ce chapitre pour y répondre.
    \item Pourquoi le positionnement CRO de BLS présenté dans ce
          chapitre (Tableau~32.2) est-il qualifié de ``raisonnement
          par analogie'' plutôt que de ``calcul''~? Quelle est la
          différence méthodologique précise~?
\end{enumerate}
\end{boxexo}

\begin{boxexo}[title={\ding{45}\quad Exercice~32.2. Checklist appliquée à un cas du cours \niveaumoyen}]
Reprenez le mécanisme de ``scellement numérique'' pour les pièces à
conviction électroniques conçu à l'Exercice~31.3
(Ch.~\ref{chap:form-meth})~: hachage SHA-256, signature ECDSA de
l'OPJ, horodatage par autorité tierce.
\begin{enumerate}[(a)]
    \item Appliquez les sept points de la checklist générale
          (Section~\ref{sec:recommandations-investigateur}) à ce
          mécanisme. Pour chaque point, indiquez si votre analyse
          antérieure (Ch.~\ref{chap:form-meth}) l'a déjà couvert,
          ou s'il reste à traiter.
    \item Ce mécanisme combine-t-il une primitive vulnérable au
          quantique (ECDSA) avec une protection post-quantique~?
          Si non, appliquez la recommandation de ce chapitre~:
          que proposeriez-vous pour sécuriser ce mécanisme à long
          terme~?
\end{enumerate}
\end{boxexo}

\begin{boxexo}[title={\ding{45}\quad Exercice~32.3. Synthèse intégrale de la Partie~VI \niveauexpert}]
Ce dernier exercice de la Partie~VI vous demande de mobiliser
\textbf{l'ensemble} des chapitres~26 à~31.

Une entreprise camerounaise de cybersécurité vous présente un
nouveau protocole, baptisé ``CamSeal'', destiné au scellement de
preuves numériques pour les tribunaux camerounais. La présentation
affirme~: \textit{``CamSeal combine un engagement de Merkle, une
preuve zk-SNARK et une signature à seuil BLS, atteignant un score
CRO de~$\mathsf{Conf}\cdot\mathsf{Rel}\cdot\mathsf{Opp} = 0{,}81$,
supérieur à tous les protocoles publiés dans la littérature.''}

\begin{enumerate}[(a)]
    \item \textbf{(Ch.~\ref{chap:CRO}-\ref{chap:prim-crypt})}~:
          évaluez la plausibilité du score annoncé en le comparant
          aux valeurs réellement publiées (Tableau~28.3,
          plage~$0{,}27$--$0{,}57$). Quelles questions précises
          poseriez-vous pour vérifier ce chiffre~?
    \item \textbf{(Ch.~\ref{chap:cryptanal})}~: appliquez la
          taxonomie des failles (Tableau~30.1) pour identifier au
          moins deux catégories de vulnérabilités à examiner en
          priorité pour cette architecture à trois couches.
    \item \textbf{(Ch.~\ref{chap:quantique}, ce chapitre)}~: la
          présentation ne mentionne aucune protection
          post-quantique. Quelle est l'implication, en particulier
          pour la couche BLS, à la lumière de la
          Section~\ref{sec:etude-bls} de ce chapitre~? Que
          recommanderiez-vous~?
    \item \textbf{(Ch.~\ref{chap:form-meth})}~: l'entreprise
          affirme avoir ``prouvé formellement'' la sécurité de
          CamSeal avec Tamarin, sans préciser si certains lemmes
          restent à l'état ``non prouvé''. Quelle question
          poseriez-vous pour clarifier ce point, et pourquoi
          est-elle essentielle avant tout déploiement judiciaire~?
    \item \textbf{(Synthèse, ce chapitre)}~: appliquez la checklist
          complète (Section~\ref{sec:recommandations-investigateur})
          à CamSeal, en indiquant pour chaque point ce que la
          présentation a déjà établi, et ce qui resterait à
          vérifier avant que vous, en tant qu'expert judiciaire
          camerounais (Ch.~12~(UE~SEC4031)), puissiez
          recommander l'adoption de ce protocole pour un usage
          judiciaire.
    \item Rédigez un court avis technique de synthèse (une demi-page),
          destiné à l'autorité qui vous a mandaté, résumant votre
          niveau de confiance actuel dans CamSeal et les conditions
          nécessaires avant tout déploiement opérationnel.
\end{enumerate}
\end{boxexo}

\bigskip

\begin{boxinfo}
\textbf{Conclusion de la Partie~VI.} Ce chapitre clôt la Partie~VI
de ce cours, consacrée aux fondements théoriques avancés~: l'impact
du quantique sur la preuve (Ch.~\ref{chap:quantique}), le Trilemme
CRO formalisé et publié (Ch.~\ref{chap:CRO}), son application
méthodique à des primitives réelles (Ch.~\ref{chap:prim-crypt}),
la lecture critique d'un protocole en prépublication
(Ch.~\ref{chap:ZKNR}), les principes de conception et de
cryptanalyse (Ch.~\ref{chap:cryptanal}), l'outillage de la preuve
formelle (Ch.~\ref{chap:form-meth}), et enfin cette mise en pratique
intégrative.

Le fil conducteur de ces sept chapitres a été, au-delà de leur
contenu technique propre, une \textbf{discipline méthodologique
constante}~: distinguer rigoureusement ce qui est établi, publié et
vérifiable, de ce qui demeure une hypothèse de travail ou une
question ouverte. Cette discipline n'est pas spécifique à la
cryptographie théorique, c'est exactement celle exigée de tout
rapport forensique depuis le Ch.~\ref{chap:rapport} de ce cours.

\medskip
\textbf{Conclusion du cours.} Ce chapitre clôt aussi le parcours
entamé au Prologue~: de l'affaire MBONGO, où une simple erreur
procédurale a anéanti un dossier judiciaire, à la question de la
pérennité cryptographique des preuves face à une rupture
technologique encore incertaine. Entre ces deux extrêmes, vous
avez acquis la rigueur méthodologique du PIU, la maîtrise
technique de l'acquisition et de l'analyse multi-supports, la
compréhension du cadre juridique camerounais et international,
et la posture éthique et critique indispensable à la pratique
de l'investigation numérique.

\textit{Cette dernière compétence, l'esprit critique face aux
affirmations non vérifiées, qu'elles proviennent d'un suspect,
d'un outil automatisé, ou d'un discours commercial sur une
technologie émergente, est sans doute la plus durable de
toutes celles enseignées dans ce cours. Les outils changeront.
Les lois évolueront. La rigueur intellectuelle, elle, reste.}
\end{boxinfo}